\section{Validator Incentive Mechanism}

\label{sec:validators}
%%==========================================================================

\subsection{Participation Decision Model}

A Primary Network validator $v$ faces a binary decision for each appchain $s$:
participate (opt-in) or abstain. The net benefit of participation is:

\[
  \Delta U_{v,s} = \underbrace{R_{v,s}}_{\text{Appchain reward}} - \underbrace{C_{v,s}^{\text{infra}}}_{\text{Additional hardware}} - \underbrace{C_{v,s}^{\text{ops}}}_{\text{Operational overhead}}.
\]

Validator $v$ opts in if $\Delta U_{v,s} > 0$.

\subsection{Opportunity Cost Constraint}

A more complete model recognizes that validator time and attention are scarce. If a
validator can commit to at most $K$ appchains simultaneously (due to hardware and bandwidth
limits), then participation in appchain $s$ may crowd out a more profitable appchain $s'$:

\[
  \Delta U_{v,s}^{\text{net}} = \Delta U_{v,s} - \max_{s' \neq s, s' \in \text{available}} \Delta U_{v,s'}.
\]

\begin{proposition}[Threshold Participation]
Under the opportunity cost model, a validator $v$ participates in the $k$ appchains with
the highest $\Delta U_{v,s}$ values, up to capacity $K$. Appchain $s$ achieves full validator
participation only if $\Delta U_{v,s} > 0$ for at least $n_s^{\min}$ validators.
\end{proposition}

This creates a competitive market among appchains for validator attention, which we analyze
in Section~\ref{sec:elastic}.

\subsection{Reward Structure Design}

Appchain $s$ can structure validator rewards in several ways. We analyze three mechanisms:

\paragraph{Mechanism A: Fixed Annual Yield.}
Appchain $s$ offers validators an annualized yield $r_s$ on their Primary Network stake weight:
\[
  R_{v,s}^A = r_s \cdot w_v \cdot P_{\text{LUX}}.
\]

\textit{Advantage}: Simple and predictable. \\
\textit{Disadvantage}: Appchain pays even when utilization is low.

\paragraph{Mechanism B: Usage-Based Rewards.}
Appchain $s$ distributes all collected fees to validators proportional to participation:
\[
  R_{v,s}^B = \frac{w_v}{\sum_{v' \in \mathcal{V}_s} w_{v'}} \cdot \text{FeeRevenue}(s, \text{epoch}).
\]

\textit{Advantage}: Aligned with appchain utilization. \\
\textit{Disadvantage}: Revenue volatility may deter participation.

\paragraph{Mechanism C: Hybrid (Recommended).}
Base fixed yield plus usage bonus:
\[
  R_{v,s}^C = \underbrace{r_s^{\text{base}} \cdot w_v \cdot P_{\text{LUX}}}_{\text{fixed base}} + \underbrace{\beta_s \cdot \frac{w_v}{\sum_{v'} w_{v'}} \cdot \text{FeeRevenue}(s, \text{epoch})}_{\text{usage bonus}},
\]
where $r_s^{\text{base}}$ is the guaranteed base yield and $\beta_s \in (0,1)$ is the
fraction of fee revenue shared with validators.

\begin{theorem}[Incentive Compatibility of Mechanism C]
Under Mechanism C, a validator $v$ with capacity $K \geq 1$ participates in appchain $s$ if:
\[
  r_s^{\text{base}} + \beta_s \cdot \frac{\mathbb{E}[\text{FeeRevenue}(s)]}{n_s \cdot \bar{w} \cdot P_{\text{LUX}}} > r_{\text{primary}} + \delta_{\min}(v),
\]
where $\delta_{\min}(v)$ is validator $v$'s minimum premium (accounting for opportunity cost).
\end{theorem}

\begin{proof}
Validator $v$ participates when $\Delta U_{v,s}^C > 0$. Substituting $R_{v,s}^C$, dividing
by $w_v \cdot P_{\text{LUX}}$, and recognizing that the expected usage bonus equals the
right-hand side expression gives the stated condition.
\end{proof}

\subsection{Empirical Calibration}

We calibrate $\delta_k$ (the minimum appchain premium by tier) against revealed preferences
from 47 production appchains. Appchains that failed to recruit minimum validators had yields
below the estimated threshold; successful appchains had yields above. Logistic regression
on appchain-validator adoption events yields:

\begin{table}[H]
\centering
\caption{Empirically Estimated Minimum Validator Premiums by Tier}
\label{tab:premiums}
\begin{tabular}{lrrr}
\toprule
\textbf{Tier} & \textbf{$\delta_k$ (median)} & \textbf{$\delta_k$ (P10)} & \textbf{$\delta_k$ (P90)} \\
\midrule
T1 & 2.1\% & 1.2\% & 3.8\% \\
T2 & 3.4\% & 2.1\% & 5.7\% \\
T3 & 4.8\% & 3.2\% & 7.2\% \\
T4 & 6.2\% & 4.5\% & 9.1\% \\
T5 & 8.7\% & 6.8\% & 12.4\% \\
\bottomrule
\end{tabular}
\end{table}

Higher-tier appchains require larger premiums because they impose greater hardware demands.
The wide P10--P90 intervals reflect heterogeneity in validator risk preferences and
existing hardware configurations.

%%==========================================================================
